\subsection{Changes of rings}

PROPOSITION 7.--- Let \(\rho: A \to B\) be a local homomorphism of Noetherian local rings, M a finitely generated A-module and N a finitely generated B-module. Put \(\overline{B} = B \otimes_A \kappa_A = B / \rho(m_A).B\) and \(\overline{N} = N \otimes_B \overline{B} = N / \rho(m_A).N\) . We have

\[
\dim_ {B} (M \otimes_ {A} N) \leqslant \dim_ {A} (M) + \dim_ {B} (\overline {{N}}),\tag{12}
\]

and there is equality if N is flat over A.

One may suppose M and N non-zero. According to the corollary of th. 1 (n° 2), there exists a subset S (resp. T) of \(m_A\) (resp. \(m_B\) ) such that \(\text{Card}(S) = \dim_A(M)\) (resp. \(\text{Card}(T) = \dim_{\overline{B}}(\overline{N})\) ) and that M/SM is an A-module of finite length (resp. \(\overline{N}/\overline{T}\overline{N}\) is a B-module of finite length). One has \(\rho(S) \subset m_B\) . Let E be the B-module \(M \otimes_A N\) ; the B-module \(E/(\rho(S).E + T.E)\) is isomorphic to M/SM \(\otimes_A N/TN\) , hence is of finite length: according to prop. 18 and 19 of II, § 4, n° 4, one has

\[
\begin{array}{r l} \operatorname{Supp} (\mathrm{M} / \mathrm{SM} \otimes_ {\mathrm{A}} \mathrm{N} / \mathrm{TN}) & = \operatorname{Supp} (\mathrm{N} / \mathrm{TN}) \cap {} ^ {a} \rho^ {- 1} (\operatorname{Supp} (\mathrm{M} / \mathrm{SM})) \\ & = \operatorname{Supp} (\mathrm{N} / \mathrm{TN}) \cap {} ^ {a} \rho^ {- 1} (\operatorname{Supp} (\kappa_ {\mathrm{A}})) \\ & = \operatorname{Supp} (\mathrm{N} / \mathrm{TN} \otimes_ {\mathrm{A}} \kappa_ {\mathrm{A}}) = \operatorname{Supp} (\overline {{\mathrm{N}}} / \mathrm{T} \overline {{\mathrm{N}}}) = \left\{\mathfrak {m} _ {\mathrm{B}} \right\}. \end{array}
\]

It follows, according to th. 1, that one has the inequality

\[
\dim_ {\mathrm{B}} (\mathrm{E}) \leqslant \operatorname{Card} (\mathrm{S}) + \operatorname{Card} (\mathrm{T}) = \dim_ {\mathrm{A}} (\mathrm{M}) + \dim_ {\overline {{\mathrm{B}}}} (\overline {{\mathrm{N}}}).
\]

Suppose now N flat over A, and let us show that there is equality in formula (12). Let a (resp. b) be the annihilator of M (resp. N). According to prop. 18 and 19 of II, § 4, n° 4, one has

\[
\begin{array}{r l} \operatorname{Supp} (E) & = \operatorname{Supp} (N) \cap {} ^ {a} \rho^ {- 1} (\operatorname{Supp} (M)) \\ & = V (b) \cap {} ^ {a} \rho^ {- 1} (V (a)) = V (b + \rho (a). B). \end{array}
\]

One consequently has

\[
\dim_ {B} (M \otimes_ {A} N) = \dim (B / (b + \rho (a). B))\tag{13}
\]

and also

\[
\dim_ {\mathbf {A}} (\mathbf {M}) + \dim_ {\overline {{\mathbf {B}}}} (\overline {{\mathbf {N}}}) = \dim (\mathrm{A} / \mathfrak {a}) + \dim (\mathbf {B} / (\mathfrak {b} + \rho (m _ {\mathbf {A}}). \mathbf {B}))  .\tag{14}
\]

Put A' = A/a and B' = B/(b + ρ(a).B) and let ρ':A' → B' be the local homomorphism deduced from ρ by passing to quotients. Since the annihilator of N is b, N is a finitely generated B/b-module with support equal to Spec(B/b), and flat over A. The homomorphism A → B/b deduced from ρ therefore possesses property (PM) of § 2, n° 1 (loc. cit., remark 3); by extension of scalars (loc. cit., prop. 1, a)), one deduces that ρ' possesses property (PM). According to loc. cit., prop. 2, one therefore has

\[
\dim (\mathbf {B} ^ {\prime}) \geqslant \dim (\mathbf {A} ^ {\prime}) + \dim (\mathbf {B} ^ {\prime} / \rho^ {\prime} (m _ {\mathbf {A} ^ {\prime}}). \mathbf {B} ^ {\prime})\tag{15}
\]

and since the ring \(\mathbf{B}'/\rho'(\mathfrak{m}_{\mathbf{A}'})\mathbf{B}'\) is isomorphic to \(\mathbf{B}/(\mathfrak{b} + \rho(\mathfrak{m}_{\mathbf{A}})\mathbf{B})\), our assertion follows from formulas (13), (14) and (15).

COROLLARY 1. --- Let \(\rho: A \to B\) be a local homomorphism of Noetherian local rings.

\begin{enumerate}
\def\labelenumi{\alph{enumi})}
\tightlist
\item
  We have
\end{enumerate}

\[
\dim (\mathbf {B}) \leqslant \dim (\mathbf {A}) + \dim (\mathbf {B} \otimes_ {\mathbf {A}} \kappa_ {\mathbf {A}})\tag{16}
\]

with equality if ρ makes B into a flat A-module.

\begin{enumerate}
\def\labelenumi{\alph{enumi})}
\setcounter{enumi}{1}
\tightlist
\item
  Suppose that \(\rho\) makes B into a flat A-module, and let S be a subset of \(\mathfrak{m}_{\mathbf{A}}\) . Then S is secant for A if and only if \(\rho(\mathbf{S})\) is secant for B.
\end{enumerate}

The assertion \(a)\) is the particular case \(\mathbf{M} = \mathbf{A}\) , \(\mathbf{N} = \mathbf{B}\) from prop. 7.

Let us prove b). Under the stated hypotheses, we have

\[
\dim (\mathbf {B}) = \dim (\mathbf {A}) + \dim (\overline {{\mathbf {B}}})\tag{17}
\]

with \(\overline{\mathbf{B}} = \mathbf{B} \otimes_{\mathbf{A}} \kappa_{\mathbf{A}} = \mathrm{B} / \rho(\mathfrak{m}_{\mathbf{A}}) . \mathbf{B}\) . Since \(\rho\) is injective (I, § 3, no. 5, prop. 9), we have

\[
\operatorname{Card} (\rho (S)) = \operatorname{Card} (S).\tag{18}
\]

Finally, B′ = B/ρ(S). B is a flat module over A′ = A/SA, hence

\[
\dim (\mathrm{B} / \rho (\mathrm{S}). \mathrm{B}) = \dim (\mathrm{A} / \mathrm{SA}) + \dim (\overline {{\mathrm{B}}})\tag{19}
\]

since \(\mathbf{B}^{\prime}/\rho(\mathfrak{m}_{\mathbf{A}^{\prime}})\mathbf{B}^{\prime}\) is isomorphic to \(\overline{B}\). Our assertion follows immediately from formulas (17), (18), and (19).

COROLLARY 2. — Let ρ : A → B be a homomorphism of Noetherian rings. Then

\[
\dim (\mathbf {B}) \leqslant \dim (\mathbf {A}) + \sup _ {\mathfrak {p} \in \operatorname{Spec} (\mathbf {A})} \dim (\mathbf {B} \otimes_ {\mathbf {A}} \kappa (\mathfrak {p})).
\]

Let q be a prime ideal of B and \(\mathfrak{p} = \rho^{-1}(\mathfrak{q})\) . By cor. 1, we have \(\dim(B_{\mathfrak{q}}) \leqslant \dim(A_{\mathfrak{p}}) + \dim(B_{\mathfrak{q}} \otimes_{A} \kappa(\mathfrak{p}))\). From this we deduce (prop. 6, b) of § 1, no. 3) the inequality \(\dim(B_{\mathfrak{q}}) \leqslant \dim(A) + \dim(B \otimes_{A} \kappa(\mathfrak{p}))\) , and we conclude by prop. 8 of § 1, no. 3.

COROLLARY 3. --- Let A be a Noetherian ring and n an integer ≥ 0. Then

\[
\dim (\mathrm{A} [ \mathrm{X} _ {1},..., \mathrm{X} _ {n} ]) = \dim (\mathrm{A}) + n.
\]

Set \(\mathbf{B} = \mathbf{A}[\mathbf{X}_1, \ldots, \mathbf{X}_n]\). For every prime ideal \(\mathfrak{p}\) of \(\mathbf{A}\) , the ring \(\mathbf{B} \otimes_{\mathbf{A}} \kappa(\mathfrak{p})\) is a polynomial ring in \(n\) variables over a field, hence has dimension \(n\) ( \(\S 2\) , no 4, cor. 1 to th. 3). By cor. 2, we have \(\dim(\mathbf{B}) \leqslant \dim(\mathbf{A}) + n\); the reverse inequality follows from Example 4 of \(\S 1\) , no. 3.

COROLLARY 4. --- Let \(\rho: A \to B\) be a homomorphism of Noetherian rings, and a an ideal of A. We have the inequality

\[
\operatorname{ht} (\rho (\mathfrak {a}). \mathbf {B}) \leqslant \operatorname{ht} (\mathfrak {a})\tag{20}
\]

if the mapping \(^a\rho : \text{Spec}(\mathbf{B}) \to \text{Spec}(\mathbf{A})\) is surjective. If B is a faithfully flat A-module, then we have \(\operatorname{ht}(\rho(\mathfrak{a}).\mathbf{B}) = \operatorname{ht}(\mathfrak{a})\) .

If B is faithfully flat over A, then \(^{a}\rho\) is surjective (II, § 2, n° 5, cor. 4 to prop. 11) and one has ht(a) ≤ ht(ρ(a).B) (§ 2, n° 1, corollary of prop. 2).

It therefore remains to prove inequality (20) under the assumption that \(^a\rho\) is surjective. Let \(\mathfrak{p}\) be a prime ideal of A such that \(\mathfrak{a} \subset \mathfrak{p}\) and \(\mathrm{ht}(\mathfrak{a}) = \mathrm{ht}(\mathfrak{p})\) ( \(\S 1, n^o 3, \text{prop. } 7\) ). Let \(\overline{\mathbf{B}} = \mathbf{B} \otimes_{\mathbf{A}} \kappa(\mathfrak{p})\) and denote by \(h\) the canonical homomorphism from B to \(\overline{\mathbf{B}}\) . If \(X = ^a\rho^{-1}(\mathfrak{p})\) is the nonempty set of prime ideals of B lying over \(\mathfrak{p}\) , we know ( \(\S 2, n^o 1, \text{lemme 1}\) ) that the map \(^ah\) is a homeomorphism from \(\operatorname{Spec}(\overline{\mathbf{B}})\) onto the subspace X of \(\operatorname{Spec}(\mathbf{B})\) . Let \(\mathfrak{q}\) be the image under \(^ah\) of a minimal prime ideal of \(\overline{\mathbf{B}}\) ; one has \(\dim(\mathbf{B}_{\mathfrak{q}} \otimes_{\mathbf{A}} \kappa(\mathfrak{p})) = 0\) and Cor. 1 implies the inequality \(\dim(\mathbf{B}_{\mathfrak{q}}) \leqslant \dim(\mathbf{A}_{\mathfrak{p}})\) . Finally,

\[
\operatorname{ht} (\rho (a). B) \leqslant \operatorname{ht} (q) = \dim (B _ {q}) \leqslant \dim (A _ {p}) = \operatorname{ht} (p) = \operatorname{ht} (a),
\]

whence the corollary.

PROPOSITION 8.--- Let A be a Noetherian ring, a an ideal of A, M a finitely generated A-module, \(\hat{\mathbf{A}}\) and \(\hat{\mathbf{M}}\) the separated completions of A and M respectively for the a-adic topology.

\begin{enumerate}
\def\labelenumi{\alph{enumi})}
\item
  Let m be a prime ideal of A containing a and \(\hat{\mathfrak{m}} = \mathfrak{m}\hat{\mathbf{A}}\) . Then \(\hat{\mathfrak{m}}\) is a prime ideal of \(\hat{\mathbf{A}}\) and one has \(\dim_{\hat{\mathbf{A}}_{\hat{\mathfrak{m}}}}(\hat{\mathbf{M}}_{\hat{\mathfrak{m}}}) = \dim_{\mathbf{A}_{\mathfrak{m}}}(\mathbf{M}_{\mathfrak{m}})\) .
\item
  We have \(\dim_{\hat{\mathbf{A}}}(\hat{\mathbf{M}}) = \sup_{\mathfrak{m}} \dim_{\mathbf{A}_{\mathfrak{m}}}(\mathbf{M}_{\mathfrak{m}})\) , where \(\mathfrak{m}\) ranges over the set of prime ideals
\end{enumerate}

(resp. maximal ideals) of A containing a. In particular, one has \(\dim_{\hat{\mathbf{A}}}(\hat{\mathbf{M}}) \leqslant \dim_{\mathbf{A}}(\mathbf{M})\) .

\begin{enumerate}
\def\labelenumi{\alph{enumi})}
\item
  Since \(\hat{\mathbf{A}}/\hat{\mathfrak{m}}\) is identified with \(\mathbf{A}/\mathfrak{m}\) , \(\hat{\mathfrak{m}}\) is a prime ideal of \(\hat{\mathbf{A}}\) . By th. 3 of III, § 3, n° 4, \(\hat{\mathbf{A}}\) is flat over \(\mathbf{A}\) , hence \(\hat{\mathbf{A}}_{\hat{\mathfrak{m}}}\) is flat over \(\mathbf{A}_{\mathfrak{m}}\) . Moreover, the canonical map from \(\mathbf{A}\) into \(\hat{\mathbf{A}}\) induces an isomorphism of \(\mathbf{A}/\mathfrak{a}\) onto \(\hat{\mathbf{A}}/\mathfrak{a}\hat{\mathbf{A}}\) , hence also an isomorphism of \(\mathbf{A}_{\mathfrak{m}}/\mathfrak{mA}_{\mathfrak{m}}\) onto \(\hat{\mathbf{A}}_{\hat{\mathfrak{m}}}/\mathfrak{m}\hat{\mathbf{A}}_{\hat{\mathfrak{m}}}\) . We conclude by applying prop. 7 to the rings \(\mathbf{A}_{\mathfrak{m}}\) and \(\hat{\mathbf{A}}_{\hat{\mathfrak{m}}}\) and to the modules \(\mathbf{M}_{\mathfrak{m}}\) and \(\hat{\mathbf{A}}_{\hat{\mathfrak{m}}}\); the module \(\mathbf{M}_{\mathfrak{m}} \otimes_{\mathbf{A}_{\mathfrak{m}}} \hat{\mathbf{A}}_{\hat{\mathfrak{m}}}\) is isomorphic to \(\hat{\mathbf{M}}_{\hat{\mathfrak{m}}}\) (III, loc. cit. and prop. 8).
\item
  By prop. 8 of III, § 3, n° 4, the map m ↦ m̂ is a bijection from the set of maximal ideals of A containing a onto the set of maximal ideals of Â. Assertion b) follows from this and from prop. 9 of § 1, n° 4.
\end{enumerate}

COROLLARY 1. --- Let A be a Zariski ring (III, § 3, n° 3, def. 2). For every finitely generated A-module M, we have \(\dim_{\hat{\mathbf{A}}}(\hat{\mathbf{M}}) = \dim_{\mathbf{A}}(\mathbf{M})\) .

Indeed, the topology of A is the a-adic topology, where a is an ideal contained in the radical of A (loc. cit.), that is, contained in every maximal ideal m of A. It therefore suffices to apply assertion b) of prop. 8.

COROLLARY 2. --- Let A be a Noetherian ring, a an ideal of A, and \(\hat{A}\) the separated completion of A for the a-adic topology. We have dim( \(\hat{A}\) ) \(\leqslant\) dim(A), with equality when A is local and \(a\) is distinct from A.

COROLLARY 3. --- Let A be a Noetherian ring and n an integer ≥ 0. Then

\[
\dim (\mathrm{A} [ [ \mathrm{X} _ {1},..., \mathrm{X} _ {n} ] ]) = \dim (\mathrm{A}) + n.
\]

The ring A{[}{[}X₁, \ldots, Xₙ{]}{]} is the separated completion of the polynomial ring A{[}X₁, \ldots, Xₙ{]} for the \(\mathfrak{a}\)-adic topology, where \(\mathfrak{a}\) is the ideal generated by X₁, \ldots, Xₙ; we therefore have

\[
\dim (\mathrm{A} [ [ \mathrm{X} _ {1},..., \mathrm{X} _ {n} ] ]) \leqslant \dim (\mathrm{A} [ \mathrm{X} _ {1},..., \mathrm{X} _ {n} ])
\]

by cor. 2. We also have

\[
\dim (\mathrm{A} [ \mathrm{X} _ {1},..., \mathrm{X} _ {n} ]) = \dim (\mathrm{A}) + n
\]

by cor. 3 of prop. 7. Finally, we have

\[
\dim (\mathrm{A}) + n \leqslant \dim (\mathrm{A} [ [ \mathrm{X} _ {1},..., \mathrm{X} _ {n} ] ])
\]

by example 4 of § 1, n° 3.

Remarks. --- 1) Let A be a Noetherian ring and a an ideal of A. Suppose A is separated and complete for the a-adic topology, and consider the ring R = A \{ X \(_{1}\) , \ldots, X \(_{n}\) \} of restricted formal series (III, § 4, n° 2, def. 2). We have dim(R) = dim(A) + n. Indeed, R is the completion of the ring B = A{[}X \(_{1}\) , \ldots, X \(_{n}\) {]} for the aB-adic topology, whence dim(R) ≤ dim(A{[}X \(_{1}\) , \ldots, X \(_{n}\) {]}) = dim(A) + n. The reverse inequality is proved as in the case of formal series.

\begin{enumerate}
\def\labelenumi{\arabic{enumi})}
\setcounter{enumi}{1}
\tightlist
\item
  Let A be a Noetherian local ring, identified with a subring of its completion \(\hat{\mathbf{A}}\) , and let B be a subring of \(\hat{\mathbf{A}}\) containing A. Suppose that B is Noetherian local and that we have \(m_{\mathbf{A}}B = m_{\mathbf{B}}\) . Then the canonical injection of A into B extends to an isomorphism of \(\hat{\mathbf{A}}\) onto the completion \(\hat{\mathbf{B}}\) of B (III, § 3, no. 5, prop. 11), whence dim(B) = dim(A) (cor. 2 to prop. 8). This applies in particular to the following situation. * Let k be a complete non-discrete valued field, A the local ring of the polynomial ring \(k[\mathrm{X}_1, ..., \mathrm{X}_n]\) at the prime ideal generated by \(\mathrm{X}_1, ..., \mathrm{X}_n\) , and B the ring of convergent power series in \(\mathrm{X}_1, ..., \mathrm{X}_n\) with coefficients in k. Then the preceding hypotheses are satisfied, and consequently one has dim(B) = n. *
\end{enumerate}

